\chapter{超越人妻}

這本書的本意，是讓讀者們都能獲得快樂且成為更優秀的自己。而與人妻約會，是現代社會中對大多數人活得快樂的最低成本方式，所以前面很多篇幅詳細講述和人妻約會的方方面面。但是在很多時候，我們和人妻約會，本質上也是為了我們自身的快樂。而快樂不一定需要在人妻之上獲得。也可以找未婚的小妹妹，甚至很多人沒意識到，快樂也不一定要從女人身上獲得。而這一章，也是回到了書的本意，告訴大家和人妻約會之外，一個優秀的男人如何自我實現，並獲得內心的快樂。

\section{避免沉船}

「沉船」是一個來自香港粵語的俗語，一開始指的是愛上妓女，後面又泛指愛上不應該愛上的人（炮友之類的）。

在和人妻約會的時候，沉船的可能性確實非常大。首先人妻因為有經驗，更懂得怎麼哄男人開心。其次又因為人妻從你身上要求的少很多，所以你會感到更多的溫柔的感覺。而且人妻更有性經驗，可能在床上會更開心。所以男人其實稍有不慎很容易在人妻身上沉船。而且一旦和人妻開始有愛情的感覺，那麼人妻的安全感需求，包括對小孩的撫養責任，就會自然而然地放在你身上。如果人妻還是凌菲菲這種邪惡的人，可能會像凌菲菲吸李松堅一樣，把所有的時間和金錢都吸乾，給她和前夫的小孩。

所以在和人妻約會遇到沉船的時候，其實需要退後一步，重新審視這段關係。一開始你尋找人妻，不就只是因為成本低嗎？現在因為人妻不要回報地對你好，你愛上了人妻。但是人妻能提供給你的也是有限的不是嗎？陪伴、生育、養老，很多事情都無法提供。當你退後一步，把這些事想清楚了，你只是想要一個廉價的替代品，那麼也就不會沉船了。

\section{回歸少女}

前面說到，人妻有很多優點。比如容易泡，又比如一般不需要出錢負責。對於剛出社會的年輕中國男人，是性價比極高的一個選擇。但是隨着男性的個人進步與發展，一個高價值的人，是需要開始遠離人妻的。

之前說過，如果結婚相當於買車的話，那麼有租車、有網約車、有拼車。而約人妻就相當於直接走進隔壁鄰居老王家把車開走。對大部分人來說，是性價比最高的開車方式。但是我現在還是沒有去開鄰居家的車，我買了一輛奔馳S級轎車，每天開車。為什麼？因為十幾萬美金對我來說跟不用錢一樣，所以我可以買一輛奔馳，什麼時候想開就開，而且也不用怕髒。對我來說的一筆小錢，讓我生活自由度和舒適度大幅提高。畢竟可能在郭菊陽老公讀到這本書的性病防治部分以前，已經自己不注意患上了HPV之類的病了吧。公用的多髒啊。

而說到第一點想開就開，就要講到當時21年我第一次回國的時候第一次聯繫上郭菊陽的事情了。當時我剛好因為準備要把賣房所得的幾千萬人民幣換成美金投資美國科技股，於是我就來到了中國灰產之都深圳找人換錢。當時恰好郭菊陽在深圳，順手發了個微信約她出來。一開始她約不出來，一下子說老公在家了，一下子說要帶小孩了。因為我當時還有事需要回北京，沒法在深圳呆太久，郭菊陽這麼浪費我時間，我當場就發怒了，於是我就直接打了個30000元的微信轉賬過去，說「你拿去僱個保姆然後出來陪我。」於是她就出來了，雖然最後她也沒收我的錢。

這說明什麼？對我來說，我的時間成本很高。跟別人共享一個人妻，對我來說我損失的時間成本，比如如果在深圳待多一兩天再去北京的成本，肯定就超過30000人民幣了。也就是說，當你是一個年輕而收入普通的國男，你去找郭菊陽或者許潔自然非常省錢，畢竟有她們老公付生活費。\textbf{但是當你提高自己，讓你的時間成本已經高到一定境界的時候，你就真正成長了。你已經超越了人妻，又可以回歸少女了。}

在23年前9個月的時候，阿爾法星還是一個一人公司\footnote{其實一開始連公司都沒有註冊}。我一個人帶着一臺手提電腦全世界邊旅遊邊創業，有時候在深圳約身份證4開頭的妹妹們，有時候在上海約身份證3開頭的妹妹們。當時我在上海，約了一個從福建客家山區去張江當初級程式設計師的小妹妹出去玩，開着李松堅的邁巴赫S680，從她在世紀大道的家裏開車帶她去上海迪士尼，拿着一大疊Priority Pass\footnote{Priority Pass是迪士尼樂園裏面景點的收費插隊的特殊門票，官方名稱是迪士尼尊享卡（Disney Premier Access）}玩一整天，然後晚上再到外灘5號的茹絲葵牛排館（Ruth's Chris Steak House）吃了一頓戰斧牛排大餐。一整天下來，不算李松堅的邁巴赫，也就花了我5000人民幣左右。從迪士尼的魔幻城堡到外灘的迷醉的霓虹燈下，看着福建妹妹在上海灘望着我迷醉的表情，這5000塊，不比那30000塊值錢很多倍？

\section{超越女性}

今年\footnote{2025年}四月，我有一天晚上在深圳南山區的某個五星級海景酒店和一個浙江女人約會完。事後那個女的就跟我作鬧起來了，我跟她吵了幾十分鐘，實在受不了，於是就直接打車回福田了。當時在計程車上，看着購物公園和Coco Park酒吧街的燈光，恍惚之間，感覺那個\textbf{女人就像《聊齋》裏面的女鬼一樣，來到我身邊，要吸取我的陽氣。如果我讓她把我陽氣都吸光之後，我就死掉了。}

盯着酒吧的燈光的時候，我又突然回過神來。轉念一想，我恍惚之間的幻想，其實難道只是幻想嗎？我交往過這麼多女性，沒有一個女性在相處中不是對我進行索取，消耗。這種索取、消耗，本質上和女鬼吸取陽氣又有什麼區別呢？

只要和女性接近過，你肯定能感受到這種消耗。這種接近不一定是要有過性關係或者結婚，只要確定男女朋友，或者稍微曖昧一下，甚至你跟你周邊女同學女同事稍微多聊幾句，你都能感受到這種消耗。你在學校有沒有經歷過關係稍微親近的女生叫你幫她寫作業？你在工作的時候是不是幫女同事搬桶裝水？你在微信加了一個女人\footnote{現在很多人用「女性」，甚至「女生」指代成年婦女。本書不會有這種語言腐敗。}之後，是不是聊兩句天，女的就要跟你要微信轉賬零花錢？之前的胖貓自殺事件，在全亞洲鬧得沸沸揚揚，起因就是胖貓被一個重慶的撈女譚竹用要零花錢的方式掏空了所有金錢然後自殺。

而一旦關係升級到男女朋友甚至夫妻關係，這種索取就會升級。比如在中國非常流行的要求工資卡上交，男性只能偷偷藏「私房錢」。這在全球任何其他國家都是聞所未聞的。甚至每年春節聯歡晚會都會有的私房錢梗，把這種敲骨吸髓式的剝削當作娛樂，實在是惡俗不堪。

除此之外，還有在中國農村很流行的彩禮和中國城市很流行的買房加名，也是現在中國最常見的索取。尤其是買房加名，很多人沒意識到買房加名的嚴重性。在一二線城市，一套房子動輒幾百萬上千萬，即使只算女人能分走的一半，也遠遠比全國最貴的江西省的彩禮要高。最近被判刑的蘇享茂案的罪犯翟欣欣，在和蘇享茂接觸之前，就是靠離婚分得了一套北京的房產，進而用這套房產去欺騙蘇享茂。而我的父親李松堅，他找的小三凌菲菲，先是靠謀殺我姐姐和找黑社會逼我媽放棄財產離婚，再靠和我父親離婚結婚騙取上海幾十套房產，包括一套價值3億以上的檀宮別墅。\footnote{詳情見\hyperref[sister]{附錄A.4《我的姐姐李新瑩》}}

我聽過很多北京、上海的朋友，喜歡歧視外地人，說什麼「只有外地人才給彩禮，我們北京人/上海人沒有彩禮這回事」。但是實際上，彩禮最多的江西，現在普遍的彩禮不超過一百萬。而北京人上海人被騙房產的數額，動輒百萬、千萬計。最極端的凌菲菲，直接騙走我的父親李松堅幾十億。一個陌生女人，以傳統婚姻為誘餌，騙取北上廣的男人領證，沒有生小孩，過一段時間，就能把一個男人的所有積蓄掏空。而且這種掠奪比彩禮更加隱蔽。一般女人會用很多話術去包裝，讓那些被掠奪的城市人不知不覺，甚至還以為是真愛，覺得女的和他離婚是自己做得不好。下面是一些女人要房產加名的時候常見的話術，大家以後見到的時候可以注意分辨：
\begin{enumerate}
\item 你不幫我買房就是不愛我。
\item 我嫁給你們家，我需要一個保障。
\item 你不給我加名是不是防着我不信任我。如果這樣那還不如不結婚算了。
\item 我們結婚了就是一家人了。寫誰的名字不都一樣。
\item 只要是你出資，加名的時候沒有區別，離婚的時候我也分不走。
\end{enumerate}

前面四點實在太過可笑，沒有必要浪費篇幅在這裏反駁。但是最後一點還是很有迷惑性的。因為一般來說，女的還會配合一兩個似是而非的案例和女權自媒體的傳銷號的宣傳。但是中國並不是判例法國家，即使有一兩個女方分不了房產的特殊案例，也代表不了你離婚的時候不會被分房產。房產加名，在全世界各個國家，基本就是默認就是贈予房產價值的50\%。如果貸款只有你一個名字而房產有女方的名字，默認就是女方有房產無貸款價值的50\%，你也有房產無貸款價值的50\%，但是房貸只屬於你一個人。據我所知，在中國以外的地方也都是如此。如果女方拿出案例告訴你有離婚判決是按照婚前出資比例分配，那麼只證明一件事：中國法院的離婚判決非常隨意。落到你頭上，可能法官往另外一個方向隨意一下，你拿到的就比默認50\%的比例還要少了。

其實比起零花錢、彩禮、房產加名這種直接掠奪更可怕的是，共同生活之後，女性對男性的無窮盡的人身、精神、經濟控制。

也是今年四月的時候，我一個好朋友和一個長得超級漂亮的女人約會。做愛完他想睡覺的時候，結果這個女人一定要他給個紅包，不然就一直吵不給睡覺。我當時在深圳，打着跨洋電話和他說到這件事，不斷感到後怕。其實只要睡一起，同個屋檐下。本質上，你的所有物理弱點就暴露了。比如這個女的，就可以通過不讓你睡覺，逼你打錢。而在那種狀況下，你根本沒有太多反制措施。越是高價值的男性，你越經不起女人的這番作鬧。在這種鬧事博弈之中，算到最後，一個高價值\footnote{衡量一個人價值的一個很好的簡單指標，就是一個人的單位時間經濟產出。}的人能做出的最理智的行為，其實也就是花錢安撫了事。而無論過錯在誰都理智花錢安撫，也是我到目前為止還沒出現在（別人寫的）PDF裏面的原因。

說到這件事，不得不說，對公司的管理其實也是一回事。如果你有一個上升期的盈利公司，跟任何普通員工糾結五險一金和離職補償都是低價值且不理智的行為。我創立的量化公司阿爾法星研究，入職過的員工也挺多了。到目前為止，我們所有辭退的員工，即使有着100\%的過錯，我都會非常慷慨地給超過法律要求的離職補償（severance package）。避免高價值的我本人和公司陷入離職員工的報復性破壞和官司之中。當然了，未來如果出現衍復投資的高亢之類的盜取公司商業機密離職開公司的事情，我自然也會和2 Sigma一樣把對方送進監獄的。

為什麼我這麼做，因為一個員工進入了你的公司，相當於給予了這個員工對公司的巨大破壞力。而安撫員工，減少這種破壞的可能性，是不得不付出成本的：無論是離職包之類的安撫成本，還是把對方送進監獄的司法成本。畢竟要成為南山必勝客，也是要付出成本上供的。一樣的，一個高價值男人，絕對不要隨隨便便讓一個女人進入你的生活。一旦一個女人進入了你的生活，她們可以破壞你人生的方式太多了，而你不得不付出巨大的代價才能回歸正常。尤其是在中國，司法體系對公司資方的保護力，絕對是遠遠強於男女關係中對高價值的一方（通常是男方）的保護。在這種文化制度下，我的公司尚且需要面試8輪，寧缺毋濫，篩選掉超過99\%的求職者\footnote{其實很多被拒的求職者，可以說是誤殺。但是在現在的文化制度下，不得不極度謹慎。}。為何你們在現代社會竟然敢隨隨便便交女朋友甚至結婚呢？

前面說的不讓睡覺，也算是一種比較低劣且極端的情況了。但是我遇到的很多女人，即使剛剛認識沒多久，即使她還腳踏兩隻船，也會故意在你家留下頭飾、絲襪之類的女性用品，對你之後的女性好友宣示主權。我記得2014年我剛到美國讀研究生的時候，有一個深二代女留學生，她男朋友是一個深圳的潮汕人。然後邀請我去她家。她聊到她男朋友，然後跟我抱怨潮汕家庭小孩多，然後她男朋友不是長子沒法繼承太多家產，她很不爽。然後說着說着，她就開始喝紅酒，然後就在我面前哭。說實話，我到今天還是莫名其妙。難道她覺得她的舊男朋友（苦主）和新男朋友（我）都是潮汕人就不算出軌？雖然她還繼續跟她男朋友在一起，但是我出於好心，有幾次我開車去紐約，還是讓她坐副駕同路。然後她就故意在我的車上留下了女性用品，試圖破壞我和別的女人的關係。最後她也成功了，造成了我和別的女人的爭吵。

而其實在對於有價值的男性來說最大的危險在於，如果與女性共處，女性就會不斷地通過作鬧消耗你，消耗你的精力和金錢。這種消耗是女性本能的行為。如果你向女人的作鬧妥協了，你就會被女人榨乾所有的時間精力。然後賺不到錢，然後女人便會判定你是低價值的人，然後離開你找下一個——這也是女人的本能行為，無法控制的。女人這種無理由作鬧，消耗男人精力的行為，只要有男人跟她接近就會觸發。就像之前章節說的，郭菊陽作為以前背叛過我的二婚帶男娃女，約出來都敢對我擺臉色作鬧。

作為一個男人，最重要的，是想清楚自己想要什麼，自己的夢想是什麼。我在我的微信簽名寫到：短期目標100億美金，中期目標眾議院議長，長期目標殖民半人馬座阿爾法星。即使你的目標沒有如此宏大，其實也可以：比如想要給父母更好的生活，比如想要老婆孩子熱炕頭。或者你的目標就是低俗到就想要多操幾個女人，多生幾個小孩都行。甚至你有一些奇怪的癖好比如你要收集完12星座女朋友或者原潮陽縣\footnote{大概是現在的汕頭市潮陽區和潮南區}內每個鎮的女朋友都行。每當你被女性控制，喪失自我的時候，想一下自己的目標和夢想是什麼。真的是像李松堅一樣幫凌菲菲養別人的小孩嗎？還是跟李松堅一樣把通過腐敗到手的金錢進貢給做過雞的許潔？

\textbf{每次想要被女性控制的時候，回去想一下你自己的夢想，就不會誤入歧途了。}

\section{兄弟會}

這一章的內容，算是本文最深的內容。從1644神州陸沉至今近400年，中國的道德體系從未徹底撥亂反正。在現代的中國，全社會厚黑學大行其道，以欺騙和背刺為榮，以仁義禮智信為恥。所以我在本文講的觀點，很多中國的讀者，應該是會反對甚至憤怒。筆者並不在意。

在現代中國社會，仁義禮智信被打為漢人和儒家的封建糟粕，而基督教的道德觀又無法引進，再加上文革前後官方思想一片混亂，於是在一片精神空虛之中，幾乎全民皈依了七八十年代美國好萊塢宣揚的舶來的「愛情神教」。尤其是女人和李松堅之類的小人，更是篤信愛情神教。我每個月至少和5個新認識的大齡剩女聊天，非常神奇，不管是中國還是美國，南方還是北方，她們好像腦子被memcpy\footnote{記憶體複製}似的說出一模一樣的話：「我還沒有遇到過有感覺的人。」而李松堅之類的小人更是離譜，為了所謂真愛，能故意讓我姐姐得上肺結核差點死掉，然後再故意讓我姐姐傳染肺結核給我\footnote{詳見本書附錄收錄文章\hyperref[sister]{《我的姐姐李新瑩》}}。幸好我抵抗力強，最後沒有被傳染。也就是說，一旦成為愛情神教的教徒，無論男女，似乎仁義禮智信都能置之腦後，甚至殺子殺女也能接受。

而愛情，說到底，就是人類促進交配繁殖的神經系統和內分泌系統的綜合反應罷了。豬豬狗狗也一樣有。所以我現在看到那些愛上我的女的，就一點波瀾也沒有——誰會不喜歡一個高中的時候就是全美國電腦奧賽第一名金牌，白手起家身價幾十億，然後拿起筆當劍，對抗不公的社會，寫下無數文章以重塑民族精神的人呢？我是女的我也流水。郭菊陽每次對我發情的心理，更像是在2009年賣飛NVIDIA的股市韭菜。

而現在，每個月都有好多撲上來找我的女人。在我眼中，她們像是拿着100塊人民幣要去買一個比特幣，然後說她認識誰誰誰當年100塊買了比特幣了。她們還會跟我說我要相信愛情，然後愛她們，然後再領證把50\%的資產甚至100\%的資產奉獻給她們，如果不這麼做就是褻瀆愛情。我每次聽一個外貌完全不同的新的女人說出一字不差的話，都覺得特別可笑。

\textbf{其實人類真正偉大的情感，是多個不存在血緣關係的男性之間，能夠有兄弟一般的情感，一起為了一個共同的目標而奮鬥。}

如果觀察歷史，幾乎所有偉大的事業，都是來源於多個男性，靠着高於純粹利益的目的，形成一個共同體完成的。這種精神，在歷朝歷代都是所有人歌頌的對象。從桃園三結義，到白帝城託孤和《出師表》。從社會頂層的文人、進士們，到社會最底層的戲劇，都對那種忠義讚嘆。從運城到成都到汕頭，人們都在祭拜、悼念關公和諸葛丞相。可惜到如今，很多人看到白帝城託孤，第一反應是劉備在試探諸葛亮。人心不古，真是讓人嘆息。

這種男性之間的兄弟般的聯盟，不僅僅是在中國傳統儒家文化中被倡導，在西方也被社會正面歌頌、倡導。我在美國多年，美國的很多文化，其實就是從小訓練人與人之間的信任與合作關係。小孩子從小就參加足球、橄欖球等重視團隊合作的項目。然後很多大學生，會加入一些學生自治組織的兄弟會（fraternity housing）。一群男大學生，自願申請加入一個不由學校或者政府領導的學生組織。一群人像兄弟一樣，在一個大房子裏同吃同住。很多人一輩子的友誼就是在兄弟會裏面建立起來的。

筆者在美國留學的時候，也曾在麻省理工學院（MIT）裏面一個類似的兄弟會裏居住過。每天住在一起，和兄弟們一起做飯、喝酒、打乒乓球。本章節的標題「兄弟會」，也是取自美國大學的兄弟會組織。

畢業後我也在美國的矽谷的科技公司工作過好長時間。矽谷有一種創業文化，就是一個有錢人，買了一個矽谷的大房子，然後讓一群有能力、有理想、但是沒有錢的年輕人住到他的房子裏面創業。屋主不收房租，免費給吃給喝，甚至還會給創業者一些資源，只需要換取一小部分的新公司股權。我在矽谷有一個朋友，結婚前就把一群參加過USACO（美國電腦奧林匹克競賽）的朋友都叫到他家的大房子住一起。現在這群同住過的人，好多都在矽谷或者創業或者投資，身價加起來已經不知道有多少億美金了。而我的公司剛開始創業的時候，我在深圳雖然沒有矽谷的五百萬美金大別墅，但是也租了一個出租屋，幾個男性朋友住在一起，吃着豬腳飯和腸粉創業。也算是把那套矽谷的創業文化因地制宜地移植了過來。

很多人低估了一個以信念和理想而不只是利益結合在一起的一群男人的戰鬥力。我在小公司和超大公司都呆過。那種幾十萬人上百萬人的大型組織，每天的權力鬥爭、繁文縟節、腐敗貪污、無意義內耗，大概率浪費了整個組織80\%以上的勞動力。其實只要很少的人，只要能減少這些內耗，就能爆發巨大的力量。

我本人其實經歷過類似的組織，是我的高中汕頭市金山中學的電腦奧林匹克競賽班（金中競賽班）。當我剛剛加入金中競賽班的時候，金中歷史上連一個進入省隊的人都沒有。雖然汕頭市第一中學有過IOI銀牌，但是汕頭市金山中學一個都沒有。高一進入金中競賽班不久，我就成為事實上的金中競賽班的領導者。這個領導者，不是學校任命的，而是金中競賽班的成員們內部公認的。作為金中競賽班成績最好的一位，我很熱心地教同級和比我小的學弟學妹們。

更重要的是，真正的領導，不僅僅是權力，更是一份責任。當時金中基本上什麼教學資源都沒有，我就用我當時還很一般的英語，跟美國電腦競賽國家隊的教練要美國隊的內部訓練資料分享給競賽班的隊友們。然後當時金山中學的校長李麗麗，對整個競賽班充滿敵意。只要哪個學生晚上去競賽班時間比較長，她就會命令所有老師煽動協同所有學生去霸凌那個競賽班的學生。當時在我高二的時候，在廣東的所謂「小高考」\footnote{當年廣東有一個考試叫做「學業水平測試」，只要及格就能順利拿到畢業證。對於金山中學這種重點中學的學生來說都很簡單。但是當時李麗麗和金山中學的老師們把這個考試叫做「小高考」，並且把它當作一個對學生進行服從性測試和壓迫的理由。}之前，政教處幾個人跑到競賽班，抓了幾個學生一起去政教處進行批鬥。當時被抓過去的三四個學生，本來是站一排的，我直接向前邁出一步，直接跟政教處老師爭論。我說這個考試，對我們的前途一點影響都沒有。我們的時間我們自己分配，我們的前途屬於自己，不用你們管。然後我直接對政教處的老師說，要不要打賭，看看我明天考試是不是全A。

當時金中競賽班的氛圍是非常友愛的，真的比現代很多親兄弟還要更親。當時我寫的網誌文章《飛廈forever》\footnote{原載我的網誌：https://sinyalee.com/blog/?p=236}裏面我想到我從飛廈中學到金山中學的學弟們，我當時寫到：「聽說中山那邊的OI水平都是由師兄帶師弟上來的。而我呢，從高二很少上飛廈的競賽班，也沒能憑藉自己的能力輔導他們，只能讓Sivon他們來讀我和炫圭艱深難懂的程序。所以我真的很內疚，覺得自己沒有盡到做師兄的責任。」

這種友愛也是互相的，我在網誌文章《For WH》\footnote{原載我的網誌：https://sinyalee.com/blog/?p=383}裏面寫到，吳宏當時對我說：「你很有天賦，OI\footnote{電腦奧林匹克競賽}可以進省隊的，可是還是要努力。還有你書也會讀好的。你的那個兩年之誓，很讓人震撼，也給了我很多感觸。說真的，NOI\footnote{全國電腦奧林匹克競賽}金牌曾經是我一個夢想，現在不可能了，我希望你可以。」

當到我高二終於能代表廣東到全國比賽之後，接觸到湖南、江蘇、上海那些全中國頂級電腦競賽名校的學生，才知道我們廣東和汕頭的這種氛圍是多麼難得。在那些（廣東的）北方的學校裏面，電腦競賽班同學之間，互相仇恨，背刺是常態。如果一個學生有了一些自己拿到的學習資料，都會藏着掖着不敢讓同學知道。當時高二第一次出省比賽，深刻地感受到了那些學生的人品和氛圍，也是我把我的大學升學志願定為第一麻省理工學院、第二香港科技大學、第三廣州中山大學的原因之一\footnote{詳見\hyperref[3faces]{附錄A.3《我的父親李松堅》}}。

這種兄弟般的團結和友愛之下，在我領導金中競賽班的那一年，汕頭金中從一個默默無聞的三線城市學校，成為廣東排名第一的電腦競賽強校。2009年的全國電腦奧林匹克競賽，全國20個金牌，汕頭就佔了2個，佔了全國的1/10。

而現在再重新回看，在我們那一屆之後，我帶起來的金中競賽班，又出了月之暗面創始人楊植麟，「腦王」鄭林楷等人。在我和我們那一屆同學的爭取和構建的風氣之下，金中競賽班繼續下去，成為一個不斷培養人才的「兄弟會」。

\begin{figure}[H]
\centering
\includegraphics[width=2.2in]{figures/gdoi.jpg}
\caption{我那一屆的電腦奧賽廣東省賽城市排名}
\end{figure}

\begin{figure}[H]
\centering
\includegraphics[width=2.5in]{figures/stoi.jpg}
\caption{廣東省賽我們獲得市第一名和學校第一名兩個獎杯}
\end{figure}

這就是忠義和男性友誼的優勢，在個體層面，可能一個體系裏面的小人可以佔盡便宜，但是如果一個群體能夠通過信仰、忠義或者友誼團結起來，超越囚徒困境，這個群體的戰鬥力，是遠遠大於一個由小人組成的群體的。而一個由信仰和忠義確立的群體，堅決排除、懲罰小人，也是絕對正義的。

\section{舔狗與小人}

上一章從正面講偉大的男性友誼和偉大的人是怎麼樣的。這一章，會從相反的角度講一下，那些小人是怎麼樣的。

在孔子和儒家的理論中，有一對非常重要的概念：「君子」和「小人」。儒家對君子和小人的定義其實很清晰：君子喻於義，小人喻於利。也就是說，君子是講道義的，而小人只講利益。以這個標準看，今天的中國，遍地都是小人，君子越來越少了。

之前說的，能夠為了理想情同兄弟的人，自然是屬於君子。小人沒有忠義，是不可能有這種男性間的友誼和同盟關係的。你想一下，你身邊是不是有很多這種人，一旦發現你優秀，他們就嫉妒你，仇恨你，落井下石，想要把你拉下水？這些就是孔子所謂的小人了。尤其是我就讀過的清華大學。無論是教授還是學生，絕對是小人濃度最高的一個地方。我在金中競賽班，即使我是成績最好的一個，大家對我的稱呼都是「新野」（Sinya）或者「新野師兄」。但是在清華，所有人一見面都是「野神」、「野牛」地叫。不只是對我這樣，人與人互相之間，就是充滿了這種假情假意的恭維。而在恭維背後，是無盡的落井下石。比如我在清華姚班上學的時候，我曾經問過同學期末考時間、課程評分標準之類的問題，結果都有同學故意告訴我錯誤的期末考時間，故意告訴我錯誤的評分標準，甚至最終導致我有一門課程掛科。

清華大學這種可怕的氛圍，我直到2022年才重新體會到。當時因為加入了城堡證券（Citadel Securities），搬到了芝加哥。這也是我除了在清華四年居住在北京之外，人生唯一一次在內陸城市居住。當時晚上7點多加班結束後，太陽已經下山，一個人步行在芝加哥的CBD洛普區（Chicago Loop）的街道上。路邊都是遊蕩的無家可歸的黑人，褲子提到屁股中間，露出半條溝子。這些人看到我，像鬣狗看到獵物一樣，惡狠狠地盯着我的眼睛，大概率是在看我什麼時候露出破綻，可以上前咬一口。我一方面不敢直視那些黑人，怕激怒他們，另一方面又不停用雙目的餘光注意那些人的一舉一動，防止他們什麼時候突然對我進行偷襲。而我在清華，感受是一模一樣的。在那一刻，芝加哥那些遊蕩的黑人，比城堡證券的同事更像我的清華同學。

說到我在清華的經歷，就不得不說清華大學兩大糟粕文化之一的「女生節」\footnote{另一大糟粕文化是「清華特獎」。}。每年3月7日就是清華大學所謂的女生節。在這個節日，清華的男生會組織起來，所有男生進行一場舔狗表演。男生要拼命貶低自己的人格去單方面舔女生。誰舔得越賣力，誰就越光榮。

當時我們2010級姚班在女生節的時候，有幾個舔狗頭子，要求我們姚班25個左右的男生，每個男生必須出幾百塊，給5個女生每人買幾千塊的禮物。當時我嚴正反對了幾句，但是還是懶得跟那些人吵架，最後還是出了幾百塊。當時大學的時候，我賬戶裏有500萬人民幣的零花錢。而我知道一些姚班的同班同學，一個月零花錢不到1000人民幣，我為他們發聲，他們卻一聲不吭，從不到一千的零花錢裏面掏出幾百給比他們富裕得多的女同學們。（我跟女生很熟，我知道她們的家庭狀況。）

清華的這些男生當舔狗，甚至是到了一種匪夷所思的地步。就是我如果對一個女的付出，肯定是想要娶這個女的或者上這個女的。是希望有回報的。但是清華這群舔狗，我觀察他們，他們似乎是不追求正面回應，甚至是在期望負面回應的痛苦的。越是被女人拒絕，他們就越享受那種痛苦的爽感。

也就是說，清華學生這個群體，對女人付出一切的舔狗和那些背刺同學的小人，都尤其多。這其實並不是巧合。\textbf{每一個舔狗都是小人，沒有例外。}一個人要成為舔狗，必然首先要放棄自身的尊嚴。一個連自身尊嚴都可以不要的人，也就是一個把利益放在一切之上的人。這些人，自然為了利益可以背信棄義。他們把女人當作物品、戰利品，然後放棄尊嚴去獲取，同時會仇視、敵對那些他們認為有競爭關係的同性，落井下石。這就是那些背刺同學的清華學生同時也是女生節的舔狗的原因。

我在《我的父親李松堅》\footnote{《我的父親李松堅》全文收錄於\hyperref[3faces]{附錄A.3}}一文中寫的李松堅的三張臉，一張是「對女神奉獻一切的臉」，還有一張是「對員工驕橫跋扈」的臉。當時寫的這兩張臉，其實不只是李松堅，而是代表一整類人：每個為女神喪失尊嚴奉獻一切的舔狗，都是會欺壓、背刺身邊男性的小人。

雖然這篇本書主要是打算寫給男性閱讀，但是這對女性也有很強的借鑑意義。對女性來說，一個男人一旦不請自來給你當舔狗，殷勤獻得太過，這個男人大概率是要盡量避開的。一個小人給你當舔狗的唯一目的就是得到你，而對小人來說，放棄尊嚴得到之後自然是不珍惜的。這樣說來，其實撈女狠狠爆舔狗金幣，還是有其正義性的。